\section{An integral comparison for the given Wach lattice}

Fix a prime $p$, the original tied period data $(P,A_{\inf})$, a good
layer $(K,s,L)$, the coefficient ring $O=\mathcal O_L$, and the given
cyclotomic group, series and inclusion records satisfying
$\mathrm{IoW}$. Fix the given \'etale functor $E$ and Wach datum
$\mathrm{Mo}$ on it. For a Galois-stable free $O$-lattice $T$ in
$V$, write $N=N(T)$ for its Wach module, $j$ for the given embedding of
its \'etale module, and $n=[K_0:\mathbb Q_p]$. Write
\[
 A=\prod_{i\in\mathbb Z/n}\!O[[X]],\qquad
 A_0=O\otimes_{\mathbb Z_p}W(\mathcal O_{\mathbb C_p^\flat}),
 \qquad W_0=O\otimes_{\mathbb Z_p}W(\mathbb C_p^\flat).
\]
The series map is $\iota:A\to A_0$, and $A_0\hookrightarrow W_0$
is the integral inclusion. Write $M_A=A_0\otimes_O T$ and
$M_W=W_0\otimes_O T$. We treat first a crystalline $V$.

% BEGIN L78W1
\paragraph{L78W1. A basis of the original lattice over the product of series rings.}
The Wach module $N$ is free over $A$ of rank $d=\operatorname{rank}_O T$.
Here is how the source's weaker projectivity over $\mathbb Z_p[[X]]$
suffices. As a direct summand of a free module over that ring, $N$
has injective multiplication by $X$. Reduction of its split
embedding modulo $X$ also shows that $N/XN$ has no $p$-torsion.
For example, if $p x=Xy$, the constant coefficient of every free
coordinate of $x$ vanishes. Divide those coordinates by $X$ and
apply the splitting to obtain $x=Xz$ in $N$.

For an idempotent $e_i$ of $A$ put $N_i=e_iN$. The finite $O$-module
$N_i/XN_i$ is torsion-free: every nonzero element of the DVR $O$
divides a power of $p$, and $p$ is injective on the quotient.
Lift an $O$-basis from this quotient. Nakayama's lemma makes the
lifts generators of $N_i$ over $O[[X]]$. They are independent:
a relation has coefficients divisible by $X$ after reduction;
injectivity of $X$ gives the divided relation. Iteration and
$X$-adic separatedness of $O[[X]]$ force all coefficients to vanish.
Let $d_i$ be this rank. The source's series inclusion into the
\'etale coefficient ring is injective, so the image of $e_i$ is a
nonzero idempotent there. On that nonzero factor the defining
base change of $N$ is the original free \'etale module of rank $d$.
The invariant basis number property gives $d_i=d$. Assemble the
component bases in the same order to obtain the asserted basis.
% END L78W1

% BEGIN L78W2
\paragraph{L78W2. Frobenius forces integral coefficients.}
Suppose a rectangular matrix $Z$ over $W_0$ satisfies
\[
             \varphi(Z)=SZS'+Y
\]
with $S,S',Y$ over $A_0$, using the original inclusions and
coefficientwise Witt Frobenius. Every entry of $Z$ lies in $A_0$.
First work without $O$ and put $F=\mathbb C_p^\flat$,
$F^+=\mathcal O_F$. The reduction of the equation modulo $p$ has
entrywise $p$th powers on the left. If the largest absolute value
$M$ of an entry of the reduced matrix were greater than one,
the corresponding left side would have absolute value $M^p$,
and every term on the right has absolute value at most $M$.
This proves integrality modulo $p$ (empty matrices are
vacuous). Choose the entrywise Teichm\"uller lift $Z_0$ of these
integral residues. Write $Z=Z_0+pZ_1$. The integral matrix
$SZ_0S'+Y-\varphi(Z_0)$ is divisible by $p$ already in $W(F^+)$:
its residue vanishes in $F$ and $F^+\hookrightarrow F$ is
injective. Here one uses perfection, so the kernel of the
residue map on Witt vectors is precisely $pW(F^+)$.
After dividing by $p$, $Z_1$ satisfies the same equation with
an integral new inhomogeneous term. Iteration gives integral
approximations modulo every $p^k$. Modulo $p^{k+1}$ the first
$k+1$ Witt coordinates agree with those of the approximation.
Thus every Witt coordinate belongs to $F^+$, proving integrality
of the actual matrix.

For coefficients in $O$, fix its finite free $\mathbb Z_p$-basis.
Represent multiplication by each tensor coefficient as a matrix over
$W(F)$. The resulting block substitution preserves multiplication,
addition, the two inclusions and Frobenius. The result just proved
applies to the enlarged rectangular matrix. For each original entry,
its multiplication matrix has coefficients in $W(F^+)$. Applying
this matrix to the integral coordinate vector of $1\in O$ recovers
the original tensor in $A_0$. This argument covers ramified $L$.
% END L78W2

% BEGIN L78W3
\paragraph{L78W3. One exponent gives both integral inclusions.}
Let $\epsilon_s$ be the $p^s$th root of the given tilt parameter,
$\pi_s=[\epsilon_s]-1$, and put $x=1\otimes\pi_1$ in $A_0$.
Then for some $m=m(T)$ one has
\[
 x^m (A_0j(N))\subset M_A,\qquad x^mM_A\subset A_0j(N)
 \quad\text{inside }M_W.
\]
Here $A_0j(N)$ denotes the $A_0$-span after applying $j$.
In particular, for each $t\in T$ the second inclusion gives
$x^m(1\otimes t)$ in the original $O$-span of the vectors
$a j(y)$ for $a\in A_0$ and $y\in N$.

To see the exponent, write $Q$ for the source's level-shifted
cyclotomic factor and $q=\iota(Q)$. In the series ring the Frobenius
sends $1+X$ to $(1+X)^p$. As $s\ge1$,
$(1+\pi_s)^{p^{s-1}}-1=\pi_1$.
The identity obtained by applying $\varphi^{s-1}$ to
$\varphi(X)=X\sum_{j=0}^{p-1}(1+X)^j$ gives $\varphi(x)=xq$;
equivalently the polynomial quotient here is the geometric sum.
Choose the basis from L78W1 and denote its images under $j$ by
$b_i$. They span $M_W$ because $N$ generates its \'etale base
change and the image of that module generates $M_W$. There are
$d$ such vectors in a free module of rank $d$ over a commutative
ring. The square surjection splits and its determinant is a unit.
They form a basis.

For every basis element the two Wach height conditions give a power
of $Q$ carrying its Frobenius into $N$, and a power carrying the
element itself into the $A$-span of $\varphi(N)$. Take $m$ to be the
maximum of all these exponents, or zero in rank zero. Semilinearity
shows that $\varphi(N)$ is contained in the span of the Frobenius
images of the chosen basis. In row coordinates relative to a fixed
$O$-basis of $T$, let $Y$ be the invertible matrix of the $b_i$.
For matrices $P,H$ over $A_0$ the two height conditions read
\[
 q^m\varphi(Y)=PY,\qquad q^mY=H\varphi(Y).
\]
Consequently $\varphi(x^mY)=P(x^mY)$. Multiplying the second
identity on the left by $Y^{-1}$ and on the right by
$\varphi(Y^{-1})$ gives
$\varphi(x^mY^{-1})=(x^mY^{-1})H$. L78W2 makes both scaled
matrices integral. The rows of the first express $x^mb_i$ in
$M_A$. Those of the second express the scaled $T$-basis in the
$A_0$-span of the very same $b_i$. Each $j(y)$ is a combination
of the $b_i$ with coefficients in $\iota(A)$. This proves both
inclusions and supplies the original $\mathrm{Sand}$ predicate.
% END L78W3

% BEGIN L78W4
\paragraph{L78W4. The period coefficient map and its clearing unit.}
The specified structural map $\iota_A:A_{\inf}\to B_{\rm cris}$
of the tied datum agrees on $\mathbb Z_p$ with the canonical map
through $\mathbb Q_p$. The two inclusions into $L\otimes_{\mathbb Q_p}
B_{\rm cris}$ thus define a ring map
\[
 \beta:A_0\longrightarrow L\otimes_{\mathbb Q_p}B_{\rm cris},
 \qquad a\otimes w\longmapsto a\otimes\iota_A(w).
\]
In particular it has the prescribed value on every simple tensor.
Tensoring the injection $A_{\inf}\to W(F)$ with the free
$\mathbb Z_p$-module $O$, and then with the free $O$-module $T$,
gives the injections used in L78W3.

The clearing factor $\beta(x)$ is a unit. Indeed, put $u=[\epsilon]$.
The compatible system of roots has untilt one and its first shifted
untilt is the prescribed nonidentity $p$th root of unity. It gives a
primitive generator of $\mathbb Z_p(1)$, so its crystalline logarithm
is a Tate period and is invertible in $B_{\rm cris}$. For clarity,
the divisibility used here is $u-1\mid p\log u$ in $A_{\rm cris}$.
In the uncompleted divided-power ring, construct $y_k$ with
$p^k(u-1)y_k=u^{p^{k+1}}-1$, starting with
$y_0=1+u+\cdots+u^{p-1}$. If $v=u^{p^{k+1}}$, then
$v-1$ is divisible by $p^{k+1}$ in that ring: for
$u=1+a$ with $a\in\ker\theta$, the terms
$\binom{p^{k+1}}j a^j$ have the factor $p^{k+1}$ after writing
$a^j=j!\,a^{[j]}$ for $j\ge1$.
Hence $1+v+\cdots+v^{p-1}=p+p^{k+1}c$ there.
Set $y_{k+1}=y_k(1+p^kc)$. These $y_k$ converge $p$-adically
and their product with $u-1$ tends to $p\log u$, by the crystalline
logarithm limit $\log u=\lim_h(u^{p^h}-1)/p^h$.
Both $p$ and the Tate period are units in $B_{\rm cris}$, hence so
is the image of $u-1$. Since $u=[\epsilon_1]^p$, the factor
$[\epsilon_1]-1$ of $u-1$ also has invertible image. This proves
the assertion for $x$ on the given layer and applies equally at $p=2$.
% END L78W4

% BEGIN L78W5
\paragraph{L78W5. The six crystalline period-map identities on the same datum.}
Put $B=B_{\rm cris}$. There is an additive map
$b_T:N\to V\otimes_{\mathbb Q_p} B$ for each $T$, together
with the given de Rham realization on $M_A$, satisfying
$\mathrm{IsPer}$ with coefficient action $\beta\iota$.
For a simple tensor $a\otimes t\in M_A$ let
$b_A(a\otimes t)=\beta(a)(t\otimes1)$, where the two coefficient
factors act on $V$ and $B$ respectively. This is $O$-balanced and
extends additively. Write $i:M_A\to M_W$ for the injection.
For crystalline $V$ and $z\in N$, L78W3 supplies $y\in M_A$ with
$i(y)=x^m j(z)$. Define
\[
                    b_T(z)=\beta(x)^{-m} b_A(y).
\]
Its more general characterization is important. Whenever
$i(y')=a'j(z)$ with $a'\in A_0$, injectivity gives
$x^m y'=a'y$. Since $\beta(x)$ is invertible,
$b_A(y')=\beta(a')b_T(z)$. Conversely the equation for
$(x^m,y)$ determines $b_T(z)$ uniquely. In particular it is
independent of the clearing exponent. For noncrystalline $V$ put
$b_T=0$, as permitted by the native contract.

Use a common clearing power to add two elements. Uniqueness gives
additivity, and multiplying an integral representative by
$\iota(f)$ gives $b_T(fz)=\beta(\iota(f))b_T(z)$ for $f\in A$.
For a Galois element $g$ or Frobenius, apply the corresponding
semilinear operation to a clearing pair $(x^m,y)$. Its first entry
still has unit image because its image is the transform of
$\beta(x^m)$. The required commutations on coefficients follow on
$O\otimes A_{\inf}$ by applying the original Galois or Witt
Frobenius map to the second tensor factor and using the ties of
$\iota_A$ and $P$. Thus $b_T(\gamma_{\rm res(g)}z)=g b_T(z)$,
and whenever $Q^k\varphi(z)\in N$,
\[
 b_T(Q^k\varphi(z))=\beta(\iota(Q^k))\,(1\otimes\varphi)b_T(z).
\]
For the latter, multiply the transformed numerator by $\iota(Q^k)$
before applying uniqueness. The same argument with a lattice
morphism $f:T\to T'$ acts as $f$ on the second tensor factor of
$M_A$ and as $f\otimes 1$ on $V\otimes B$. For crystalline source
and target, and whenever the image is in $N(T')$, this proves the
required naturality.

Finally set $\beta_{\rm dR}=(1\otimes\iota_{\rm dR,B})b_A$, where
$\iota_{\rm dR,B}:B\to B_{\rm dR}$ is the specified period-ring map.
It takes $(a\otimes w)\otimes t$ to $at\otimes\iota_{\rm dR}(w)$.
For every $k\ge0$ and $y'\in M_A$ with $i(y')=x^k j(z)$ the general
characterization above gives
\[
 \beta_{\rm dR}(y')=\iota_{\rm dR}(\pi_1)^k
                 (1\otimes\iota_{\rm dR,B})b_T(z).
\]
This is the last field of the original $\mathrm{IsPer}$.
It is an exact equality for every clearing witness, including a
power smaller than the uniform exponent in L78W3.
% END L78W5

% BEGIN L78W6
\paragraph{L78W6. Convergent evaluation on the crystalline coefficient algebra.}
On the given layer there is a single $J$ satisfying the eight fields of
$\mathrm{IsEval}$ for $\beta$ of L78W4. Its value on an open-disc
series is obtained before passing to $B_{\rm cris}$. Put
$D_0=O\otimes_{\mathbb Z_p}A_{\rm cris}$ and $D=D_0[1/p]$.
It has the norm defined by the lattice $D_0$:
$\|z\|_D=p^{-h}$ where $h$ is the largest integer with
$z\in p^hD_0$, and $\|0\|_D=0$. The integral crystalline ring is
$p$-torsion-free, separated and complete, and $O$ is finite free
over $\mathbb Z_p$. Thus the norm is nonarchimedean and
submultiplicative, and $D$ is complete: a Cauchy tail belongs to
one translate of a power of $p$ times $D_0$. Write
$c:L\to D$, $b:A_0\to D_0$, and
$\kappa:D\to L\otimes B_{\rm cris}$ for the canonical maps, using
the prescribed tie to $P$. They satisfy $\kappa b=\beta$.
One has $\|c(v)\|_D\le p|v|$ by scaling $v$ into $O$ by a power
of $p$. Frobenius and the Galois maps preserve $D_0$ and fix $c(L)$,
so act continuously on $D$.

Write $u_s=b(1\otimes\pi_s)$. In $D_0/p$ one has
$u_s^{p^s}=u_0$ and $u_0^p=0$, since $[\epsilon]-1\in\ker\theta$
has integral divided powers. Consequently, with $d=p^{s+1}$,
$\|u_s^k\|_D\le p^{-\lfloor k/d\rfloor}$. For a series $h$ on the
open unit disc, $\sum_k|h_k|r^k<\infty$ at every real $0<r<1$:
choose $r<r'<1$, use boundedness of $|h_k|(r')^k$, then sum the
geometric majorant. Thus $h(u_s)=\sum_k c(h_k)u_s^k$ exists
absolutely. Put $f=b\iota$, $d_i=f(e_i)$, and define
\[
                J(h)=\kappa\!\left(\sum_i d_i h_i(u_s)\right).
\]
Outside the componentwise convergent algebra one may define $J$
arbitrarily. The $d_i$ are orthogonal idempotents summing to one.
Absolute convergence gives the sum and Cauchy product laws and
the scalar law. For integral $h$, the evaluation agrees with
$f(h)$: the remainder after truncation is $X^m$ times an integral
series whose image under $f$ has norm at most one. This proves
$J(h)=\beta\iota(h)$ with the original, prescribed inclusion.

For Frobenius the parameter substitution is $(1+X)^p-1$ and
the idempotent $d_{i-1}$ goes to $d_i$. The Galois substitution is
$(1+X)^{\chi(g)}-1$ and fixes the idempotents. These assertions
hold for $f$ on integral series by the given inclusion and its
Witt Frobenius and Galois ties (LWR3). An integral substitution
with zero constant term preserves each disc of radius $r<1$ and
is bounded on its power series with the Gauss norm there.
Approximate the outer series by polynomials at a radius larger
than $p^{-1/d}$. Both composed evaluations are continuous by the
bound on $u_s^k$, so the integral identities extend to every
convergent series with coefficients in $L$. This proves the two
action laws with the prescribed Frobenius shift. Finally the
defining crystalline logarithm gives $\log(1+u_0)=t$ in $D$.
The identity $p^s\log(1+X)=\log((1+X)^{p^s})$ of open-disc
series, evaluated at $u_s$, shows that the constant tuple of
logarithms has image $1\otimes p^{-s}t_B$.

Here are details for the eighth field, which compares the
Banach limit with a de Rham congruence. For a fixed $m\ge1$, the
map $A_{\rm cris}\to B_{\rm dR}^+/\mathrm{Fil}^m$ is bounded on
its integral lattice for the usual $p$-adic quotient norm. In fact,
write $\ker\theta=(\xi)$ in $A_{\inf}$.
The divided powers $\xi^j/j!$ with $j\ge m$ vanish in the
quotient. For products of the others of total degree $j<m$, the
denominator divides $j!$ by the multinomial formula. Their
images thus lie in $(m!)^{-1}(A_{\inf}/\xi^m)$.
This lattice is $p$-adically complete and separated: its finite
filtration by powers of $\xi$ has successive quotients
$A_{\inf}/\xi\simeq\mathcal O_{\mathbb C_p}$, which are complete,
separated, and $p$-torsion-free. The sequences remain exact modulo
$p^k$ by torsion-freeness, and the surjective transition maps give
the assertion on taking inverse limits. Hence the bounded map on
the divided-power envelope extends to its completion. The prescribed crystalline--de Rham map agrees on the dense
divided-power envelope by its \(\mathrm{on0}\) field. Its
\(\mathrm{bdd}\) field bounds the entire integral crystalline
image in the same finite quotient by one power of $p$, so it is
continuous there and equals that extension after transport by the
ties. Under any specified
$\iota':L\to B_{\rm dR}$, the coefficients lie in
$B_{\rm dR}^+$, since an algebraic element over $\mathbb Q_p$
is integral for its filtration valuation. A finite $\mathbb Z_p$-basis
of $O$ bounds all its images modulo $\mathrm{Fil}^m$ by one further
power of $p$, using
$B_{\rm dR}^+/\mathrm{Fil}^m=(A_{\inf}/\xi^m)[1/p]$.
Thus $D\to B_{\rm dR}^+/\mathrm{Fil}^m$ is continuous. All these
bounds are for the fixed $m$ and embedding, as required.

Take $i,\zeta,\iota'$ satisfying the native agreements on $K_0$
and the chosen $p^s$th root. The map to $B_{\rm dR}$ selects $d_i$.
Indeed $d_j$ cuts out the character $\mathrm{emb}_j$ on $K_0$:
this follows on prime-to-$p$ Teichm\"uller roots from the defining
integral inclusion, componentwise multiplication by $e_j$, and
L78W4. These roots span $K_0$ over $\mathbb Q_p$ by lifting a basis
of the residue field and using reduction and Nakayama's lemma.
The distinct embeddings and $\sum_j d_j=1$ now prove the assertion.
Put $t_s=p^{-s}\iota_{\rm dR,B}(t_B)$. Modulo
$\mathrm{Fil}^m$ the image of $1+u_s$ equals
$\iota'(\zeta)\exp(t_s)$. Indeed their $p^s$th powers agree by the
nilpotent logarithm/exponential identity for $[\epsilon]$, and
their residues agree by the specified root. Their ratio is
unipotent in a nilpotent ideal of a $\mathbb Q_p$-algebra, where
raising to the $p^s$th power is injective by taking the finite
logarithm. Evaluating a polynomial therefore gives the first $m$
terms of its Taylor series in the substitution $X=\zeta\exp(Y)-1$.
For a general convergent tuple truncate $h_i$. The left side
converges by the preceding continuity, and the first $m$ coefficients
on the right converge by the definition of the given expansion
$e$. Consequently
\[
 (\iota'\otimes\iota_{\rm dR,B})_{\rm mult}J(h)
 -\sum_{k<m}\iota'(e_k)t_s^k\in\mathrm{Fil}^m B_{\rm dR}.
\]
For $m=0$ both the image and the empty sum lie in $B_{\rm dR}^+$.
This gives the asserted $J$ on all the prescribed levels and embeddings.
% END L78W6

% BEGIN L78W7
\paragraph{L78W7. The original open-disc algebra and all its expansions.}
For the same $J$ let $C$ be the product of the algebras of series over
$L$ converging on the open unit disc. If $\zeta$ has $p$-power order
set $z=\zeta$; otherwise set $z=1$, as prescribed by the source.
Define the expansion of $h$ in component $i$ by
\[
       E_{i,\zeta}(h)(Y)=h_i(z\exp Y-1)
\]
as a coefficientwise limit. This interpretation matters because the
constant coefficient of the substitution can be nonzero. Put
$y=z\exp Y-1$ and fix an output degree $k$. We have $|y_0|<1$
because a $p$-power root reduces to one. Choose
$|y_0|<r<1$ and $H\ge1$ so that $|y_j|\le rH^j$ for
$0\le j\le k$. The ultrametric inequality in the coefficient
convolution gives $|[Y^j]y^b|\le r^bH^j$ for all these $j$.
Thus the coefficient series is absolutely summable, and Cauchy
products prove that $E_{i,\zeta}$ is an $L$-algebra map.
For integral series it equals the prescribed $R.\mathrm{lev}$:
the latter sends $X$ to $z\exp(Y)-1$ (its normalized exponential
has value and derivative satisfying $E(0)=1$, $E'=E$), reads just
component $i$, and its coefficient continuity gives the same
limit of polynomial readings. This also proves agreement for the
default value $z=1$. Constants read in their indicated component.

The tuple $t(X)=\log(1+X)$ has zero constant and belongs to $C$,
since $|1/k|_p\le k$. It expands to $Y$: at a $p$-power root its
constant is $\log z=0$ by $z^{p^a}=1$, and its derivative after
substitution is $1$. Differentiation here holds coefficientwise by
the same bound at degree $k+1$, applied also to the differentiated
logarithm. At $z=1$ a factor $X^w$ gives a factor $Y^w$ because
$\exp Y-1$ has order one. The Frobenius substitution preserves
convergence, and its expansion at $(i,\zeta)$ is the expansion at
$(i-1,\zeta^p)$ with $Y$ replaced by $pY$, for the prescribed
$p$-power roots. For an allowed cyclotomic substitution the
analogous expression uses $(i,\zeta^a)$ and $\chi(g)Y$, with
$a\equiv\chi(g)\pmod{p^j}$ when $\zeta$ has order $p^j$.
Indeed the identity
$(z\exp Y)^{\chi(g)}=z^a\exp(\chi(g)Y)$ holds first for
nonnegative integers of this residue class. Take integers tending
to $\chi(g)$; their integral binomial coefficients converge
coefficientwise. The bound above controls their infinite tails
uniformly, so the formula follows degree by degree. For an
arbitrary outer $h_i$, polynomial approximation transfers the
formula: the integral inner substitutions have Gauss norm at most
the radius on any smaller disc, and each output coefficient is
continuous there at any radius strictly greater than $|z-1|$.
For the expanded side at $z^a$, choose the radius larger than
$|z^a-1|$ as well.

Now use $J$ restricted to $C$, with distinguished element $t$ and
crystalline value $p^{-s}t_B$. This value is a unit, and its de Rham
image belongs to filtration one. The element $\beta(\iota(X))$ is
a unit by factoring $[\epsilon]-1$ as a multiple of
$[\epsilon_s]-1$, as in L78W4. The character of $K_0$ on each
idempotent and the Taylor congruences are those proved in L78W6.
This gives every field of $\mathrm{Disc}$, the inclusion of every
componentwise convergent tuple, and the prescribed logarithmic
series. In particular the same $\beta$ serves for this algebra and
the lattice period maps of L78W5.
% END L78W7

% BEGIN L78W8
\paragraph{L78W8. The filtered error for a Laurent expansion in a component.}
We give the full truncation estimate used by the local comparison.
Let $C_{\iota'}$ be the submodule generated by
$c v-\iota'(c)v$ for $c\in L$ in $V\otimes_{\mathbb Q_p}B_{\rm dR}$,
using the two indicated actions. It is stable under multiplication
on the period-ring factor, as can be seen on each generator.
Denote reduction modulo it by a bar. Let $t$ be a unit in
$\mathrm{Fil}^1$, $h\in B_{\rm dR}$ arbitrary, and suppose that
the native input consists of $r$ vectors $D_j$, scalars $b_j$,
series $e_j$, an element $z$ and a Laurent sequence $L_k$ with
\[
 t^w\bar z=\sum_j b_j\bar D_j,\qquad
 b_j-\sum_{q<N}\iota'(e_{j,q})t^q\in\mathrm{Fil}^N,\qquad
 \bar L_k=\sum_j e_{j,k+w}\bar D_j,
\]
where $e_{j,q}=0$ for $q<0$. Fix $a\ge0$ with $L_k=0$ for
$k<-a$. There is an integer $c$, fixed before the truncation order
$M$, such that
\[
  h t^a z-h\sum_{k=0}^{M-1}t^k L_{k-a}
          \ \in\ C_{\iota'}+\mathrm{Fil}^{M-c}
            (V\otimes_{\mathbb Q_p}B_{\rm dR})\qquad(M\ge0).
\]
In fact choose $d$ with every $D_j\in\mathrm{Fil}^d$ by
exhaustiveness on each of the finitely many pure tensors in those
vectors. Choose $b$ with $h t^{a-w}\in\mathrm{Fil}^b$. These choices
use exhaustiveness only, including when $a<w$. Empty families
allow any $d$. For $M+w\ge a$ truncate the $b_j$ at $N=M+w-a$.
Multiplication by $h t^{a-w}$ and $D_j$ puts the error in
$\mathrm{Fil}^{M+w-a+b+d}$ by multiplicativity. In the quotient
the retained terms have the form $ht^{k+a}\bar L_k$ for
$-w\le k<M-a$. Replacing the lower limit by $-a$ preserves the
sum: $L_k=0$ below $-a$, and $\bar L_k=0$ below $-w$ by the
coefficient formula. Reindexing gives the displayed estimate
with $c=a-w-b-d$. If $M+w<a$ there are only finitely many missed
orders. For each choose a filtration level containing its error
by exhaustiveness, and increase $c$ to meet those finitely many
bounds. This supplies one integer for every $M$, with the original
component submodule retained.
% END L78W8

% BEGIN L78W9
\paragraph{L78W9. The normalized invariant frame on a small disc.}
Use column coordinates in the basis of L78W1. Let $G_g$ be the
matrix of $g$ applied to the basis, and $P$ the matrix of
$Q^m\varphi$ applied to it, where $m$ is chosen uniformly on the
basis. There is an integral matrix $P'$ with $PP'=Q^{m'}I$:
write $Q^h$ times each basis element in the integral span of the
Frobenius images, as in L78W3, and take $m'=m+h$. Galois
acts trivially modulo $X$, so $G_g(0)=I$. Semilinearity gives
$G_{gh}=G_g g(G_h)$ and commuting the two actions gives
\[
           g(Q^m)P\varphi(G_g)=Q^m G_g\,g(P).
\]
The cyclotomic action fixes the coefficient field and acts
on the formal variable by the substitution
$X\mapsto(1+X)^{\chi(g)}-1$. Indeed Frobenius is pinned by
$X\mapsto(1+X)^p-1$, and a substitution commuting with it
with zero constant and linear term $u$ must give multiplication by
$u$ in the formal coordinate $Y=\log(1+X)$: a series commuting
with $Y\mapsto pY$ has no terms of degree $k\ne1$.
Goodness, surjectivity of the cyclotomic quotient and the source's
generator valuation give $u=\chi(\gamma)\in\mathbb Z_p^\times$
with $v_p(u-1)=s$, where $s\ge2$ if $p=2$.

On component $i$ solve $G_{\gamma,i}\gamma(U_i)=U_i$ with
$U_i(0)=I$. In degree $k>0$ the unknown $k$th coefficient
appears with multiplier $1-u^k$. All remaining coefficients on
the right are integral linear combinations of coefficients of
smaller degree, since $G_{\gamma,i}$ and the substitution are
integral. On $1+p^s\mathbb Z_p$, binomial expansion gives
$v_p(u^k-1)=s+v_p(k)$: on raising to $p$ the valuation increases
by one and on raising to an exponent prime to $p$ it is unchanged
(the given dyadic range excludes cancellation). Thus division is
possible and the solution for any chosen constant column is unique.
Its coefficients, for constant vectors of norm at most one, satisfy
\[
       \|U_{i,k}\|\le p^{sk}|k!|_p^{-1}
                \le p^{k(s+1/(p-1))}.
\]
Indeed the preceding bounds are monotone and the successive ratios
are $|1-u^k|_p^{-1}$. The inverse $V_i=U_i^{-1}$ exists formally.
Its transpose solves the same kind of fixed-column equation with
the integral matrix $G_{\gamma,i}^{-\mathsf T}$, which also has
constant $I$. Thus the same bound applies to $V_i$. All components
converge at least for $|X|<p^{-(s+1/(p-1))}$. Uniqueness of columns
also gives $G_{g,i}g(U_i)=U_i$ for every $g$: this product satisfies
the $\gamma$ equation by commutativity of the group and the cocycle
identity, and has constant $I$.

The shifted Frobenius now takes a constant matrix form.
$Q(0)=p$ is a unit over $L$, so
$W_i=Q^{-m}P_i\varphi(U_{i-1})$ is a formal matrix.
The commutation identity above makes each column of $W_i$
$\gamma$-invariant. Hence $W_i=U_i C_i$ with $C_i=W_i(0)$
over $L$. It is invertible since $PP'=Q^{m'}I$, and $U_i$ and
$\varphi(U_{i-1})$ are invertible. This gives the required
constant Frobenius matrix before the analytic continuation.
% END L78W9

% BEGIN L78W10
\paragraph{L78W10. One logarithmic power continues the frame and its inverse.}
Write $t(X)=\log(1+X)$, $Y_s=(1+X)^{p^{s-1}}-1$, so
$\varphi(Y_s)=Y_sQ$. There is a scalar $\tau$ convergent everywhere
on the open unit disc with $Y_s\tau=t$ and
$Q\varphi(\tau)=p\tau$. Explicitly,
$\tau=p^{1-s}\varphi^{s-1}(t/X)$; the identity follows by
$\varphi(t)=pt$, and $t/X$ converges on the disc.

Put $w=m+m'$. The Frobenius relation from L78W9 and its
inverse-transpose give equations of the single form
\[
        Q^w Z_i=A_i\varphi(Z_{i-1})B_i
\]
for $Z=U$ and $Z=V^{\mathsf T}$. In the first case
$A_i=Q^{m'}P_i$, $B_i=C_i^{-1}$; in the second,
$A_i=Q^{2m}(P'_i)^{\mathsf T}$, $B_i=C_i^{\mathsf T}$.
For the latter identity invert
$P_i\varphi(U_{i-1})=Q^m U_i C_i$ and use
$P_i^{-1}=Q^{-m'}P'_i$. Thus $A_i$ has integral series entries
and $B_i$ is constant in each case. Set $Z_i^*=\tau^w Z_i$.
The equation becomes $Z_i^*=p^{-w}A_i\varphi(Z_{i-1}^*)B_i$:
the denominator $Q$ disappears. For $0<r<1$ the substitution
$F=(1+X)^p-1$ has Gauss norm at radius $r$ at most
$\max(p^{-1}r,r^p)$. Iterating sends that radius into the
initial disc of L78W9 after finitely many substitutions, since
each step is bounded by $r\max(p^{-1},r^{p-1})$ with the
initial $r$ an upper bound for all later radii. Iterating the
last matrix equation the same number of times expresses every
$Z_i^*$ by products of convergent integral matrices, constants,
and the final small-disc series evaluated at the iterated
substitution. For any smaller radius this proves convergence of
the coefficients with geometric decay, by bounding at $r$.
Hence $t^wU_i=Y_s^w Z_i^*$ and likewise $t^w V_i$ converge
on the entire open disc. L78W7 places these exact entries in
$C$. The original basis, $U,V$, the Galois, Frobenius, and
constant matrices and exponents $m,w$ now give every field of
the source's $\mathrm{Frame}$ on the supplied crystalline lattice.
% END L78W10

% BEGIN L78W11
\paragraph{L78W11. Laurent coefficients with their crystalline and functorial transports.}
Use the maps of L78W5--L78W7. Extend $J$ to $C[1/t]$, which
is possible since $J(t)=1\otimes p^{-s}t_B$ is invertible.
Extend $E_{i,\zeta}$ to the same localization with values in
$L((Y))$, sending $t$ to $Y$. Thus $U,V$ may be evaluated by
both maps, by L78W10. Write $v(z)$ for the column of integral
coordinates of $z\in N$ and $B_T$ for the row of the period
vectors $b_T(n_j)$ of its basis. Set
\[
 d=B_T J(U),\qquad d_{i,j}=J(e_i)d_j,\qquad
 H_{i,\zeta}(z;Y)=
       \sum_j E_{i,\zeta}\big((Vv(z))_j\big)\,d_{i,j}.
\]
The coefficient of $Y^k$ in $H$ is exactly the source's
$\mathrm{levB}$: multiplying each entry of $Vv(z)$ by $t^w$
shifts the coefficient to $k+w$, and evaluating $U$ divides
$J(t^wU)$ by $J(t)^w$. These are the original denominators.
The action and evaluation formulas extend through localization.
By L78W5 the Galois transform of the row $B_T$ is
$B_T J(G_g)$. Since $G_g g(U)=U$, every $d_j$ is fixed by
the original absolute Galois group. The $e_i$ are fixed as well,
and $K_0$ acts on $d_{i,j}$ by $\mathrm{emb}_i$ as in L78W6.
Thus every coefficient belongs to the prescribed crystalline
subspace for that embedding.

All coefficients below $-w$ vanish. At $\zeta=1$ even all
negative coefficients vanish: the formal entries
$t^w Vv(z)$ are divisible by $X^w$, and their expansions are
divisible by $Y^w$ by L78W7. Multiplication by $a\in A$ multiplies
$H_{i,\zeta}$ by $R.\mathrm{lev}_{i,\zeta}(a)$; the coefficient
convolution is finite at each index, even at a pole. Under $g$
one has $VG_g=g(V)$. The integral coordinates transform by
$G_g g(v(z))$, so for the guarded roots of L78W7 the output is
$H_{i,\zeta^a}(z;\chi(g)Y)$; its $k$th coefficient is multiplied
by $\chi(g)^k$, for all integer $k$.

For Frobenius work first in the fraction fields of the formal
series components. For any exponent $h$ with
$y=Q^h\varphi_N(z)\in N$, the relation of L78W9 gives
\[
           Vv(y)=Q^h C\,\varphi(Vv(z)).
\]
This identity can be read in $C[1/t]$: multiply through by
powers of $t$ there, or by the extra common power $Q^m$ in
the matrix relation and cancel it in each formal series domain.
Thus no evaluation at a root requires division by the value of
$Q$. On the period side $J(Q)$ is a unit: from
$Y_sQ=\varphi(Y_s)$ and $\varphi^{s-1}(X)=Y_s$ both $J(Y_s)$
and $J(\varphi(Y_s))$ are units. The first follows by applying
Frobenius to the unit $J(X)$, and the second similarly.
The period map identity on basis elements gives
$\varphi(d)=d J(C)$; the idempotent $J(e_{i-1})$ is carried to
$J(e_i)$. Consequently the coefficient on component $i$ of
$H_{i,\zeta}(y;Y)$ is
\[
 [Y^k]H_{i,\zeta}(y;Y)
  =\sum_{j\ge0}[Y^j]R.\mathrm{lev}_{i,\zeta}(Q^h)\,
     p^{\,k-j}\varphi\!\left([Y^{k-j}]H_{i-1,\zeta^p}(z;Y)\right).
\]
This proof works for the original, arbitrary $h$, which can be
smaller than the clearing power for the chosen basis.

For a lattice morphism between crystalline objects write $M$
for its matrix in their integral Wach bases; it exists because
the given functor preserves these lattices. Then
$M G_g=G'_g g(M)$. It follows that $V'MU$ is fixed by $\gamma$
in each formal component. Such series are constant: the first
nonconstant coefficient at $k>0$ would be multiplied by
$u^k\ne1$. Writing $V'MU=C_f$ gives the identities
$V'v(fz)=C_fVv(z)$ and $(f\otimes 1)d=d'J(C_f)$.
Multiplying by $J(e_i)$ reads the $i$th component of the
constant matrix. The two identities then show
$H_{i,\zeta}(fz;Y)=(f\otimes1)H_{i,\zeta}(z;Y)$ coefficientwise.
In particular the choices of normalized frames on the two
objects can be made independently.
% END L78W11

% BEGIN L78W12
\paragraph{L78W12. The complete comparison on every original good layer.}
Define $\mathrm{levB}$ by L78W11 on crystalline objects and by zero
elsewhere. Compose with $1\otimes\iota_{\rm dR,B}$ to obtain
$\mathrm{levN}$. For each $z$ let $\mathrm{cmp}(z)$ be the sum
over $i$ of the constant coefficients of $H_{i,1}(z;Y)$, taking
the same zero convention. All maps are additive.
Multiplication by $e_i$ selects just its own summand; hence
$\mathrm{cmp}(e_iz)=\mathrm{levB}_{i,1}(z)_0$ lies in the
crystalline $i$th component, and the sum lies in $D_{\rm cris}$.
Constants act $O$-linearly, and $X$ kills $\mathrm{cmp}$ because
$E_{i,1}(X)=\exp Y-1$ has zero constant and there is no pole at one.
All support, scalar and naturality formulas pass to $\mathrm{levN}$
by finite coefficient sums. L78W11 gives the two required action
formulas for $\mathrm{levB}$, including at negative indices.
The pure tensor formula for $\beta_{\rm dR}$ is L78W5.

It remains to compare these expansions with an arbitrary integral
clearing witness on the component singled out by $\iota'$ and
$\iota'(\zeta)$. Put $a_j=(t^w Vv(z))_j\in C$.
Since $UV=I$, we have
$J(t)^w b_T(z)=\sum_j J(a_j)d_j$.
In the quotient by $C_{\iota'}$ of L78W8, scalars in
$L\otimes B$ act through multiplication after the chosen two
embeddings. This follows on pure tensors directly from the
generators of $C_{\iota'}$. By the orthogonal $K_0$ eigenspace
calculation in L78W6 the compatible embedding sends
$J(e_i)$ to $1$ and all other idempotents to $0$; therefore the
last display gives the first relation of L78W8 with
$D_j=(1\otimes\iota_{\rm dR,B})d_j$,
$z$ there equal to $(1\otimes\iota_{\rm dR,B})b_T(z)$, and
$b_j$ the image of $J(a_j)$ under the two embeddings.
The Taylor field of L78W6 gives its second relation with
$e_j=E_{i,\zeta}(a_j)$, and the coefficient formula of L78W11
gives its third. For every clearing witness $i(y)=x^k j(z)$ use
$h=\iota_{\rm dR}(\pi_1)^k$ in L78W8 and substitute the last
identity of L78W5. For every allowed pole bound $a$ this gives
the required $c$ before every truncation $M$. This proves
$\mathrm{levN\_dR}$ and completes $\mathrm{CmpS}(\mathrm{Mo})$.
The choice of frame required only that $V$ be crystalline;
the maps outside this predicate still have the stipulated domains.
% END L78W12

% BEGIN L78W13
\paragraph{L78W13. The compatible Wach functor and comparison at every prime.}
Under the source's Bhatt--Scholze integral input, every original
datum $D$ of type $D7$ admits compatible
$H:\mathrm{HWOf}(D.\mathrm{loc})$,
$F:\mathrm{FonOf}(D.\mathrm{loc},D.\mathrm{locA},H)$
and $\mathrm{CmpOf}(D.\mathrm{loc},D.\mathrm{locA},H,F)$.
Indeed $D.\mathrm{locT}$ supplies the ties at every prime and
DFD12 gives crystalline realization with the finite coefficients.
At each good layer choose the group, series and inclusion records
of LWR1--LWR3. On this very inclusion FET25 constructs the \'etale
functor; WPR30 and FET26 yield its Wach structure. The integral
Herr classes on that functor are LWR5, and L78W12 supplies its
comparison. Choose this complete tuple once for the layer. The
source's three projections of that tuple give the Wach, functor,
and comparison fields with identical underlying data. The two
additional uniformizer fields follow in $B_{\rm dR}^+$:
$[\epsilon]-1\equiv\log[\epsilon]\pmod{\mathrm{Fil}^2}$, and
$[\epsilon_1]-1$ has nonzero residue $\zeta_p-1$, so its inverse
belongs to filtration zero. Choosing the tuple at each prime
gives exactly the three records displayed above. In particular
this construction requires no transition maps between separately
indexed good layers.
% END L78W13
